% =====================================================================
% Effect-size calculation — Methods subsection
% Mate-choice copying meta-analysis
%
% Mirrors scripts/1_effect_size_calculation_pipeline/calculate_effect_sizes.py
%
% Requires in preamble:
%   \usepackage{amsmath}
%   \usepackage{hyperref}   % for \autoref / \hyperref
% Citation commands assume natbib/biblatex (\citep). Citation keys:
%   borensteinIntroductionMetaanalysis2011, chinnSimpleMethodConverting2000
%
% Each numbered equation carries a \label{eq:...} so it can be referenced
% in text, e.g. \autoref{eq:hedges-g-means} or
% \hyperref[eq:hedges-g-means]{(\autoref{eq:hedges-g-means} pooled-SD form)}.
% =====================================================================

\subsection{Effect-size calculation}

We expressed every effect as both a bias-corrected standardised mean
difference (Hedges' $g$) and a (log) odds ratio, computing whichever could
be derived directly from the reported statistics and obtaining the other by
conversion (\autoref{eq:or-to-d}--\autoref{eq:g-to-or}). All formulae and
their sampling variances follow the standard meta-analytic treatments of
\citet{borensteinIntroductionMetaanalysis2011}.

\subsubsection{Hedges' $g$ from group means and standard deviations}

When a study reported the mean, standard deviation (\textit{SD}) and sample
size of a treatment and a control group, we computed the standardised mean
difference $d$ using the pooled standard deviation $S_{p}$
(\autoref{eq:pooled-sd}) and applied the small-sample bias correction $J$
(\autoref{eq:J-means}) to obtain Hedges' $g$ (\autoref{eq:hedges-g-means})
and its sampling variance (\autoref{eq:var-g-means})
\citep{borensteinIntroductionMetaanalysis2011}:

\begin{align}
  S_{p}      &= \sqrt{\frac{(n_{1}-1)\,SD_{1}^{2} + (n_{2}-1)\,SD_{2}^{2}}{n_{1}+n_{2}-2}}, \label{eq:pooled-sd} \\[4pt]
  d          &= \frac{\bar{X}_{\text{treat}} - \bar{X}_{\text{ctrl}}}{S_{p}}, \label{eq:d-means} \\[4pt]
  J          &= 1 - \frac{3}{4(n_{1}+n_{2}-2) - 1}, \label{eq:J-means} \\[4pt]
  g          &= J \, d, \label{eq:hedges-g-means} \\[4pt]
  \mathrm{Var}(g) &= \frac{n_{1}+n_{2}}{n_{1}\,n_{2}} + \frac{g^{2}}{2(n_{1}+n_{2}-2)}, \label{eq:var-g-means}
\end{align}

\noindent where $n_{1}$ and $n_{2}$ are the treatment and control sample
sizes. Where only a standard error (\textit{SE}) was reported, the standard
deviation was recovered as $SD = SE\sqrt{n}$.

\subsubsection{Hedges' $g$ from test statistics}

For studies reporting a $t$-statistic rather than group descriptives, we
converted $t$ to $d$ following \citet{borensteinIntroductionMetaanalysis2011},
distinguishing the experimental design by the degrees of freedom. For a
one-sample or paired design ($\mathit{df} = n - 1$), $d$ was obtained from
\autoref{eq:d-t-paired} and the variance from \autoref{eq:var-g-t-paired}:

\begin{align}
  d          &= \frac{t}{\sqrt{n}}, \label{eq:d-t-paired} \\[4pt]
  J          &= 1 - \frac{3}{4(n-1) - 1}, \label{eq:J-t-paired} \\[4pt]
  g          &= J \, d, \label{eq:hedges-g-t-paired} \\[4pt]
  \mathrm{Var}(g) &= \frac{1}{n} + \frac{g^{2}}{2(n-1)}. \label{eq:var-g-t-paired}
\end{align}

\noindent For an independent two-sample design ($\mathit{df} = n_{1}+n_{2}-2$),
$d$ was obtained from \autoref{eq:d-t-indep}:

\begin{align}
  d          &= t \sqrt{\frac{1}{n_{1}} + \frac{1}{n_{2}}}, \label{eq:d-t-indep} \\[4pt]
  J          &= 1 - \frac{3}{4(n_{1}+n_{2}-2) - 1}, \label{eq:J-t-indep} \\[4pt]
  g          &= J \, d, \label{eq:hedges-g-t-indep} \\[4pt]
  \mathrm{Var}(g) &= \frac{n_{1}+n_{2}}{n_{1}\,n_{2}} + \frac{g^{2}}{2(n_{1}+n_{2}-2)}. \label{eq:var-g-t-indep}
\end{align}

\noindent When individual group sizes were not reported for an
independent-samples design, we assumed an equal split, $n_{1} = n_{2} = N/2$.
Effects reported only as $F$-statistics, $\chi^{2}$ values or generalised
linear mixed-model outputs were not directly convertible and were excluded
from this step.

\subsubsection{Odds ratio from a $2\times2$ table}

For studies reporting binary choice frequencies, we summarised the outcome
as an odds ratio (\textit{OR}; \autoref{eq:or}), with the variance of the log
odds ratio obtained by the delta-method approximation
(\autoref{eq:var-log-or}) \citep{borensteinIntroductionMetaanalysis2011}:

\begin{align}
  \mathit{OR}              &= \frac{A \, D}{B \, C}, \label{eq:or} \\[4pt]
  \mathrm{Var}(\ln \mathit{OR}) &= \frac{1}{A} + \frac{1}{B} + \frac{1}{C} + \frac{1}{D}, \label{eq:var-log-or}
\end{align}

\noindent where $A$ and $B$ are the number of treatment-group individuals
choosing the option associated with the social information and the
alternative, respectively, and $C$ and $D$ the corresponding counts for the
control group. When no separate control group existed, we used the null
expectation that individuals chose each option with equal probability,
setting $C = D = n/2$.

\subsubsection{Conversion between effect-size metrics}

To place all studies on a common scale, effects available in only one metric
were converted to the other using the logistic-to-normal approximation
\citep{chinnSimpleMethodConverting2000, borensteinIntroductionMetaanalysis2011}.
A log odds ratio was converted to a standardised mean difference
(\autoref{eq:or-to-d}) and then bias-corrected to $g$
(\autoref{eq:or-to-g}):

\begin{align}
  d               &= \ln(\mathit{OR}) \times \frac{\sqrt{3}}{\pi}, \label{eq:or-to-d} \\[4pt]
  \mathrm{Var}(d) &= \mathrm{Var}(\ln \mathit{OR}) \times \frac{3}{\pi^{2}}, \label{eq:or-to-var-d} \\[4pt]
  g               &= J \, d, \qquad \mathrm{Var}(g) \approx J^{2}\,\mathrm{Var}(d), \label{eq:or-to-g}
\end{align}

\noindent with $J$ evaluated as above ($J \approx 1$ where $n$ was unknown).
Conversely, Hedges' $g$ was converted to a log odds ratio using
\autoref{eq:g-to-or}:

\begin{align}
  \ln(\mathit{OR})              &= g \times \frac{\pi}{\sqrt{3}}, \label{eq:g-to-or} \\[4pt]
  \mathrm{Var}(\ln \mathit{OR}) &= \mathrm{Var}(g) \times \frac{\pi^{2}}{3}. \label{eq:g-to-var-or}
\end{align}

\noindent These conversions are approximate and assume that the underlying
continuous trait follows a logistic distribution
\citep{chinnSimpleMethodConverting2000}; converted values were flagged in the
data set so they could be distinguished from directly calculated effects in
sensitivity analyses.
